% 章节 05：1960 原刊数学笔误（K-01 至 K-23）与重排本错字（L-01 至 L-31）校勘记
\section{校勘记}\label{sec:errata}

本中文校订版以 R. E. Kalman 1960 年发表于 \emph{Transactions of the ASME, Journal of Basic Engineering}（第 82 卷 D 辑，第 35–45 页，对应底本第 1–11 页）的原刊文本及 2002 年重排转录电子底本为依据，对全文公式 \eqref{eq:1}–\eqref{eq:44}，定理~\ref{thm:1}–\ref{thm:5}，评注~\ref{rem:a}–\ref{rem:q}，表~\ref{tab:1}，算例~\ref{ex:1}–\ref{ex:2} 以及图~\ref{fig:1}–\ref{fig:4} 进行了逐项符号推导验算与框图信号流核查。原刊实质数学或记号笔误（共 23 条，编号 \hyperref[sec:k01]{K-01} 至 \hyperref[sec:k23]{K-23}）已在中文正文中直接订正，现将各条原貌，订正结果及严格数学证明或反例备考如下；底本重排录入与光学字符识别错字（共 31 条，编号 \hyperref[sec:l-all]{L-01 至 L-31}）则分类汇总于后。本节引述原式时统一折算为本仓库标准数学记号，对纯字形特征以文字说明。

\subsection{K-01：第 2 页右栏末样本取值记号多余时间自变量}\label{sec:k01}

\begin{itemize}
\item 定位：第 2 页右栏末行，紧接第 3 页首行式~\eqref{eq:1} 之前。
\item 原貌：将随机变量 $x_1(t_1)$ 的实数取值阈值印为带正体时间自变量的 $\xi_1(\mathrm{t})$。
\item 订正：删去多余时间自变量，改为标量 $\xi_1$。
\item 依据：紧接其后的式~\eqref{eq:1} 明确写作 $\mathbb{P}\bigl(x_1(t_1) \leq \xi_1 \mid y(t_0) = \eta(t_0), \dots, y(t) = \eta(t)\bigr) = F(\xi_1)$，此处刻画单一时刻随机变量 $x_1(t_1)$ 的取值阈值，应为无时间自变量的实标量 $\xi_1$（即便视作样本函数在待估时刻的取值，对应时标亦应为斜体 $t_1$ 而非正体 $\mathrm{t}$）。
\end{itemize}

\subsection{K-02：第 3 页定理 1 条件 (A) 条件分布对称中心误写为原点}\label{sec:k02}

\begin{itemize}
\item 定位：第 3 页左栏，定理~\ref{thm:1} 条件 (\emph{A})。
\item 原貌：文字声明条件分布函数 $F(\xi)$ 关于条件均值 $\bar{\xi}$ 对称，但对称方程印为 $F(\xi - \bar{\xi}) = 1 - F(\bar{\xi} - \xi)$。
\item 订正：对连续条件分布订正为 $F(\xi) = 1 - F(2\bar{\xi} - \xi)$（若分布含离散点质量，右端取左极限 $F\bigl((2\bar{\xi} - \xi)^{-}\bigr)$）。
\item 依据：在原印方程中令偏移量 $u = \xi - \bar{\xi}$，则 $\bar{\xi} - \xi = -u$，原式对一切实数 $u$ 化为 $F(u) = 1 - F(-u)$；特别取 $u = 0$ 得 $F(0) = 1/2$，这迫使分布关于坐标原点 $0$ 对称而非关于条件均值 $\bar{\xi}$ 对称。在非退化估计问题中，给定非零观测后的条件均值 $\bar{\xi} = x_1^{\text{✫}}(t_1 \mid t) \neq 0$，结合条件 (\emph{B}) 以 $\xi \leq \bar{\xi}$ 表示原始取值坐标可知，关于均值 $\bar{\xi}$ 对称的正确关系必为 $F(\bar{\xi} + u) = 1 - F(\bar{\xi} - u)$，即 $F(\xi) = 1 - F(2\bar{\xi} - \xi)$。
\end{itemize}

\subsection{K-03：第 3 页评注 (c) 将长方观测矩阵笼统称为奇异矩阵}\label{sec:k03}

\begin{itemize}
\item 定位：第 3 页右栏，评注~\ref{rem:c} 第一段。
\item 原貌：将部分状态可观测系统 $\bm{y}(t) = \bm{M}\bm{x}(t)$ 中的矩阵 $\bm{M}$ 笼统表述为“奇异矩阵”。
\item 订正：改述为“列秩小于 $n$ 的矩阵（当补零扩充为 $n \times n$ 方阵时为奇异矩阵）”。
\item 依据：问题~\ref{prob:I} 中状态 $\bm{x}(t)$ 为 $n$ 维向量，观测 $\bm{y}(t)$ 为 $p$ 维向量（$p < n$），故 $\bm{M}$ 一般是 $p \times n$ 长方矩阵；行列式与奇异性仅对方阵有定义，无法由单拍无噪观测唯一还原全部状态分量的充要条件是零空间非平凡，即 $\operatorname{rank}(\bm{M}) < n$。
\end{itemize}

\subsection{K-04：第 3 页评注 (c) 多维推广误用联合累积分布函数在负象限的凸性}\label{sec:k04}

\begin{itemize}
\item 定位：第 3 页右栏，评注~\ref{rem:c} 末段。
\item 原貌：声称只要 $n$ 维联合条件累积分布函数 $F(\xi_1, \dots, \xi_n)$ 关于 $\xi_1 - \bar{\xi}_1, \dots, \xi_n - \bar{\xi}_n$ 对称，且在这些变量全为负的区域上为凸函数，定理~\ref{thm:1} 对向量范数损失 $\mathbb{E}\bigl[L(\|\bm{x}(t_1) - \bm{X}(t_1)\|)\bigr]$ 依然成立。
\item 订正：订正为要求条件概率分布关于条件均值向量 $\bar{\bm{\xi}}$ 中心对称，且当存在条件概率密度 $p(\bm{\xi})$ 时，该密度的所有上水平集 $\{\bm{\xi} \in \mathbb{R}^n : p(\bm{\xi}) \geq c\}$ 均为凸集（即中心对称拟凹单峰分布，且损失函数关于原点中心对称并具凸下水平集）。
\item 反例与依据：一维情形下 $F(\xi)$ 在 $\xi \leq \bar{\xi}$ 上的凸性与密度单调不减（单峰）等价，但当维数 $n \geq 2$ 时该等价性彻底失效。取最简单的二元独立标准正态分布 $F(\xi_1, \xi_2) = \Phi_0(\xi_1)\Phi_0(\xi_2)$（其中 $\Phi_0$ 与 $p_0$ 分别为一维标准正态累积分布函数与密度函数，显然满足推论的正态性假设且 $\bar{\xi}_1 = \bar{\xi}_2 = 0$），计算其在负象限内点 $(\xi_1, \xi_2) = (-0.1, -0.1)$ 处的二阶偏导矩阵行列式：
\begin{equation*}
\det\bigl(\nabla^2 F(\xi_1, \xi_2)\bigr)\big|_{(-0.1, -0.1)} = \Phi_0(-0.1)^2 p_0(-0.1)^2 \left[0.01 - \frac{p_0(-0.1)^2}{\Phi_0(-0.1)^2}\right] \approx -0.0245 < 0
\end{equation*}
由于二阶偏导矩阵行列式严格为负，必存在严格负特征值，故二元标准正态累积分布函数在负象限根本不是凸函数！Sherman\cite{r25} 所引中心对称单峰积分引理在多维空间中的正确条件是条件密度函数的上水平集为关于 $\bar{\bm{\xi}}$ 中心对称的凸集。
\end{itemize}

\subsection{K-05：第 4 页定理 2 证明段落标签与定理条件顺序颠倒}\label{sec:k05}

\begin{itemize}
\item 定位：第 4 页右栏，式~\eqref{eq:11} 下方说明句及定理~\ref{thm:2} 的证明。
\item 原貌：式~\eqref{eq:11} 后称正交投影记号受定理小写分支 (\emph{b})（指正态过程下正交投影恒等于条件数学期望）启发；紧接证明中标签 (\emph{A}) 引用式~\eqref{eq:10} 证明二次损失线性估计情形，标签 (\emph{B}) 利用定理~\ref{thm:5} (\emph{B}) 证明正态过程情形。
\item 订正：将式~\eqref{eq:11} 后的分支引用改为 (\emph{A})，并将证明中的两段标签复位以严格对齐定理~\ref{thm:2} 的条件 (\emph{A})（正态随机过程）与条件 (\emph{B})（线性估计且二次损失），同时指明条件 (\emph{A}) 隐含损失函数满足式~\eqref{eq:2}（及其多维对称单调凸水平集推广）。
\item 依据：定理~\ref{thm:2} 正文明确将“正态随机过程”列为 (\emph{A})，将“限制为线性估计且 $L(\varepsilon) = \varepsilon^2$”列为 (\emph{B})；原刊证明将两分支标号完全颠倒。
\end{itemize}

\subsection{K-06：第 5 页式 (12) 下方观测矩阵维数行列印反}\label{sec:k06}

\begin{itemize}
\item 定位：第 5 页左栏，式~\eqref{eq:12} 下方参数维数说明段。
\item 原貌：称输出矩阵 $\bm{M}(t)$ 为 $n \times p$ 矩阵。
\item 订正：改为 $p \times n$ 矩阵。
\item 依据：在输出方程 $\bm{y}(t) = \bm{M}(t)\bm{x}(t)$ 中，状态 $\bm{x}(t)$ 为 $n \times 1$ 列向量，输出 $\bm{y}(t)$ 为 $p \times 1$ 列向量，矩阵乘法维数相容性要求 $\bm{M}(t)$ 必为 $p \times n$ 矩阵（与第 6 页问题~\ref{prob:I} 定义一致）。
\end{itemize}

\subsection{K-07：第 5 页式 (15) 下方随机激励协方差矩阵误用孤立记号}\label{sec:k07}

\begin{itemize}
\item 定位：第 5 页右栏式~\eqref{eq:15} 后统计性质第三式。
\item 原貌：定义 $\mathbb{E}[\bm{u}(t)\bm{u}^{\mathrm{T}}(t)] = \bm{G}(t)$，但紧接第 6 页首行协方差级数及全文其余各处均直接写作 $\bm{Q}(t)$。
\item 订正：统一订正为 $\mathbb{E}[\bm{u}(t)\bm{u}^{\mathrm{T}}(t)] = \bm{Q}(t)$。
\item 依据：第 2 页\hyperref[sec:notation]{符号表}，第 6 页协方差级数，第 7 页式~\eqref{eq:24}，\eqref{eq:30}，\eqref{eq:32} 以及第 8 页表~\ref{tab:1} 均以 $\bm{Q}(t)$ 表示随机激励协方差矩阵，$\bm{G}(t)$ 系手稿早期残留异号。
\end{itemize}

\subsection{K-08：第 6 页问题 I 下方悬空引用不存在的定理 2-a}\label{sec:k08}

\begin{itemize}
\item 定位：第 6 页左栏，问题~\ref{prob:I} 定义后第二段首句。
\item 原貌：写作“由定理 2-\emph{a} 可知”。
\item 订正：改为“由定理~\ref{thm:2}（结合评注~\ref{rem:f} 的向量推广）可知”。
\item 依据：第 3 页虽有定理~\ref{thm:1a}，但第 4 页仅有定理~\ref{thm:2} 及评注~\ref{rem:d}–\ref{rem:g}，全文并不存在编号为 2-\emph{a} 的定理。
\end{itemize}

\subsection{K-09：第 6 页式 (18) 第二行状态转移矩阵时间自变量笔误}\label{sec:k09}

\begin{itemize}
\item 定位：第 6 页右栏，式~\eqref{eq:18} 第二行首项。
\item 原貌：对一般待估时刻 $t_1$ 印为 $\bm{\Phi}(t + 1; t)\bm{x}^{\text{✫}}(t_1 - 1 \mid t - 1)$。
\item 订正：将状态转移矩阵时标订正为 $\bm{\Phi}(t_1; t_1 - 1)\bm{x}^{\text{✫}}(t_1 - 1 \mid t - 1)$。
\item 依据：正文声明该行系将状态方程~\eqref{eq:16} 代入上一行首项 $\mathscr{P}_{\mathcal{Y}(t-1)}[\bm{x}(t_1)]$ 所得。在时刻 $t_1 - 1$ 应用式~\eqref{eq:16} 有 $\bm{x}(t_1) = \bm{\Phi}(t_1; t_1 - 1)\bm{x}(t_1 - 1) + \bm{u}(t_1 - 1)$，向子空间 $\mathcal{Y}(t - 1)$ 作正交投影后首项必为 $\bm{\Phi}(t_1; t_1 - 1)\bm{x}^{\text{✫}}(t_1 - 1 \mid t - 1)$；对时变系统，仅在特例 $t_1 = t + 1$ 下它才等于 $\bm{\Phi}(t + 1; t)$。
\end{itemize}

\subsection{K-10：第 6 页图 3 单位延迟器前后信号标签倒置及多步预测支路接点错位}\label{sec:k10}

\begin{itemize}
\item 定位：第 6 页右栏，图~\ref{fig:3}（线性最优滤波器的矩阵框图）。
\item 原貌：原图将第二求和器出口（单位延迟器左端入口）标为 $\bm{x}^{\text{✫}}(t \mid t - 1)$，将单位延迟器右端出口标为 $\bm{x}^{\text{✫}}(t + 1 \mid t)$，并将上方通往多步预测模块 $\bm{\Phi}(t + s; t + 1)$ 的支路接在单位延迟器右端出口上。
\item 订正：执行两步联动订正——(1) 互换延迟器前后信号标签，使第二求和器出口（延迟器入口）标为 $\bm{x}^{\text{✫}}(t + 1 \mid t)$，延迟器出口标为 $\bm{x}^{\text{✫}}(t \mid t - 1)$；(2) 将上方 $\bm{\Phi}(t + s; t + 1)$ 多步预测支路的引出点前移至第二求和器出口与单位延迟器入口之间的 $\bm{x}^{\text{✫}}(t + 1 \mid t)$ 节点处。
\item 依据：对照图~\ref{fig:2} 与图~\ref{fig:4} 的统一离散时间基准，单位延迟器入口恒为 $t + 1$ 时刻量，出口恒为 $t$ 时刻量。按式~\eqref{eq:20}，右端输出矩阵 $\bm{M}(t)$ 必须接收 $\bm{x}^{\text{✫}}(t \mid t - 1)$ 方能输出 $\bar{\bm{y}}(t \mid t - 1) = \bm{M}(t)\bm{x}^{\text{✫}}(t \mid t - 1)$ 并与当前观测 $\bm{y}(t)$ 相减生成新息 $\tilde{\bm{y}}(t \mid t - 1)$；同时下方反馈矩阵 $\bm{\Phi}(t + 1; t)$ 接收 $\bm{x}^{\text{✫}}(t \mid t - 1)$ 方能输出图中所标的 $\bm{x}^{\text{✫}}(t + 1 \mid t - 1)$，与新息修正项 $\bm{\Delta}^{\text{✫}}(t)\tilde{\bm{y}}(t \mid t - 1)$ 在第二求和器相加后恰生成式~\eqref{eq:21} 左端的 $\bm{x}^{\text{✫}}(t + 1 \mid t)$。进而由式~\eqref{eq:31} 首行 $\bm{x}^{\text{✫}}(t + s \mid t) = \bm{\Phi}(t + s; t + 1)\bm{x}^{\text{✫}}(t + 1 \mid t)$ 可知，上方预测算子 $\bm{\Phi}(t + s; t + 1)$ 的输入信号必须是延迟器之前的 $\bm{x}^{\text{✫}}(t + 1 \mid t)$（若接于延迟器之后不仅时标失配，而且会漏掉当拍观测 $\bm{y}(t)$ 的新息修正项）。
\end{itemize}

\subsection{K-11：第 7 页式 (25) 推导中随机激励误带条件自变量及单步余投影记号说明}\label{sec:k11}

\begin{itemize}
\item 定位：第 7 页左栏，推导式~\eqref{eq:25} 的正交关系链。
\item 原貌：在展开 $\tilde{\bm{x}}(t + 1 \mid t - 1)$ 时将过程随机激励误印为带条件时标的 $\bm{u}(t \mid t - 1)$；同段开头以残损括号式表示 $\bm{x}(t + 1) - \mathscr{P}_{\mathcal{Z}(t)}[\bm{x}(t + 1)]$。
\item 订正：将误带条件时标的激励项还原为 $\bm{u}(t)$；将单步新息子空间上的余投影规范记为 $\mathscr{P}_{\mathcal{Z}(t)}^{\perp}[\bm{x}(t + 1)] = \bm{x}(t + 1) - \mathscr{P}_{\mathcal{Z}(t)}[\bm{x}(t + 1)]$。
\item 依据：由式~\eqref{eq:16} 与 $\mathscr{P}_{\mathcal{Y}(t-1)}[\bm{u}(t)] = \bm{0}$ 可得 $\tilde{\bm{x}}(t + 1 \mid t - 1) = \bm{\Phi}(t + 1; t)\tilde{\bm{x}}(t \mid t - 1) + \bm{u}(t)$，激励向量 $\bm{u}(t)$ 本身并无条件记法，原印系受相邻项 $\tilde{\bm{x}}(t \mid t - 1)$ 视觉同化所致。
\end{itemize}

\subsection{K-12：第 7 页式 (25) 前观测标量分量个数下标误作 m}\label{sec:k12}

\begin{itemize}
\item 定位：第 7 页左栏，式~\eqref{eq:25} 正上方关于 $\bm{M}(t)$ 行线性无关的说明句。
\item 原貌：将观测向量的各个标量分量列举为 $y_1(t), \dots, y_m(t)$。
\item 订正：将末项下标由 $m$ 订正为 $p$，即 $y_1(t), \dots, y_p(t)$。
\item 依据：第 6 页问题~\ref{prob:I} 已明确规定观测向量 $\bm{y}(t)$ 为 $p$ 维向量，$\bm{M}(t)$ 为 $p \times n$ 矩阵，其行数为 $p$（字母 $m$ 在第 5 页与第 8 页中专指控制输入向量的维数）。
\end{itemize}

\subsection{K-13：第 7 页多步预测状态展开求和通项中转移矩阵右时标漏加 1}\label{sec:k13}

\begin{itemize}
\item 定位：第 7 页左栏末，对 $t_1 = t + s \geq t + 1$ 反复应用式~\eqref{eq:16} 的展开式。
\item 原貌：求和项通项印为 $\bm{\Phi}(t + s; t + r)\bm{u}(t + r)$（且求和下限漏印等号，见\hyperref[sec:l-all]{L-18}）。
\item 订正：将转移矩阵右自变量订正为 $t + r + 1$，即 $\sum_{r=1}^{s-1} \bm{\Phi}(t + s; t + r + 1)\bm{u}(t + r)$（约定 $s = 1$ 时空和为零向量 $\bm{0}$）。
\item 依据：取 $s = 2$ 检验，由式~\eqref{eq:16} 得 $\bm{x}(t + 2) = \bm{\Phi}(t + 2; t + 1)\bm{x}(t + 1) + \bm{u}(t + 1)$，对应 $r = 1$ 项系数矩阵必为单位矩阵 $\bm{I} = \bm{\Phi}(t + 2; t + 2)$，而非原式给出的 $\bm{\Phi}(t + 2; t + 1)$。物理上 $\bm{u}(t + r)$ 在时刻 $t + r + 1$ 方进入状态 $\bm{x}(t + r + 1)$，此后至时刻 $t + s$ 经历的状态转移必为 $\bm{\Phi}(t + s; t + r + 1)$。因该求和各项均与 $\mathcal{Y}(t)$ 正交，向 $\mathcal{Y}(t)$ 投影后归零，故未波及式~\eqref{eq:31}。
\end{itemize}

\subsection{K-14：第 7 页平滑情形 s < 0 时待估时标符号写反且混称滤波器阶数为向量坐标数}\label{sec:k14}

\begin{itemize}
\item 定位：第 7 页右栏定理~\ref{thm:3} 正上方一段。
\item 原貌：讨论 $s < 0$（数据平滑）时将待估量印为 $\bm{x}^{\text{✫}}(t - s \mid t)$，并称该向量具有 $(1 - s)(n - p)$ 个坐标。
\item 订正：将待估量订正为 $\bm{x}^{\text{✫}}(t + s \mid t)$，并将后半句明确表述为“生成固定滞后平滑估计 $\bm{x}^{\text{✫}}(t + s \mid t)$ 的降阶动态滤波器需要 $(1 - s)(n - p)$ 个状态变量”。
\item 依据：前文已固定记待估时刻为 $t_1 = t + s$；当 $s < 0$ 时，$t + s < t$ 恰表示滞后 $|s| = -s$ 步的平滑，而写成 $t - s = t + |s| > t$ 则反转回超前预测（后半句维数因子 $1 - s = 1 + |s|$ 正是按负数 $s < 0$ 计入滞后步数）。同时 $\bm{x}^{\text{✫}}(t + s \mid t) \in \mathbb{R}^n$ 本身恒有 $n$ 个坐标分量，$(1 - s)(n - p)$ 实为将单步 $n - p$ 维降阶滤波器（见评注~\ref{rem:j}）扩充为滞后 $-s$ 步平滑器所需的内部状态维数。
\end{itemize}

\subsection{K-15：第 7 页式 (32) 二次协方差递推方程矩阵乘积顺序颠倒}\label{sec:k15}

\begin{itemize}
\item 定位：第 7 页右栏，评注~\ref{rem:h} 式~\eqref{eq:32}（2002 年底本第 12 页转录附记已核实 1960 年原刊确印如此）。
\item 原貌：花括号内第二项末尾两个矩阵因子印为 $\bm{P}^{\text{✫}}(t)\bm{M}(t)$。
\item 订正：颠倒乘积顺序订正为 $\bm{M}(t)\bm{P}^{\text{✫}}(t)$。
\item 依据：将式~\eqref{eq:28} 与式~\eqref{eq:29} 代入式~\eqref{eq:30} 右端首项展开：
\begin{equation*}
\begin{aligned}
&\bm{\Phi}^{\text{✫}}(t + 1; t)\bm{P}^{\text{✫}}(t)\bm{\Phi}^{\mathrm{T}}(t + 1; t) = \bigl[\bm{\Phi}(t + 1; t) - \bm{\Delta}^{\text{✫}}(t)\bm{M}(t)\bigr]\bm{P}^{\text{✫}}(t)\bm{\Phi}^{\mathrm{T}}(t + 1; t) \\
= &\bm{\Phi}(t + 1; t)\bigl\{\bm{P}^{\text{✫}}(t) - \bm{P}^{\text{✫}}(t)\bm{M}^{\mathrm{T}}(t)\bigl[\bm{M}(t)\bm{P}^{\text{✫}}(t)\bm{M}^{\mathrm{T}}(t)\bigr]^{-1}\bm{M}(t)\bm{P}^{\text{✫}}(t)\bigr\}\bm{\Phi}^{\mathrm{T}}(t + 1; t)
\end{aligned}
\end{equation*}
在维数上，$\bigl[\bm{M}(t)\bm{P}^{\text{✫}}(t)\bm{M}^{\mathrm{T}}(t)\bigr]^{-1}$ 为 $p \times p$ 矩阵，其右侧必须左乘 $p \times n$ 矩阵 $\bm{M}(t)$ 再乘 $n \times n$ 矩阵 $\bm{P}^{\text{✫}}(t)$；原刊所印 $(n \times n)(p \times n)$ 在 $p < n$ 时矩阵乘法根本无定义，且破坏了误差协方差矩阵的对称性。
\end{itemize}

\subsection{K-16：第 7 页评注 (h) 反求最优滤波器系数时误引公式编号 (29–30)}\label{sec:k16}

\begin{itemize}
\item 定位：第 7 页右栏，评注~\ref{rem:h} 第二段首句。
\item 原貌：称由式~\eqref{eq:32} 算出 $\bm{P}^{\text{✫}}(t)$ 后，可立即由“公式~(\ref{eq:29}–\ref{eq:30})”确定线性最优滤波器。
\item 订正：订正为“公式 \eqref{eq:28}–\eqref{eq:29}”。
\item 依据：确定最优滤波器式~\eqref{eq:21} 所需的两个系数矩阵是最优新息加权矩阵 $\bm{\Delta}^{\text{✫}}(t)$ 与闭环状态转移矩阵 $\bm{\Phi}^{\text{✫}}(t + 1; t)$，分别由式~\eqref{eq:28} 与式~\eqref{eq:29} 给出；式~\eqref{eq:30} 则是已被消元合并进式~\eqref{eq:32} 的协方差更新方程。
\end{itemize}

\subsection{K-17：第 8 页定理 4 对偶代换规则漏写对所得各矩阵方程两边取整体转置}\label{sec:k17}

\begin{itemize}
\item 定位：第 8 页右栏，定理~\ref{thm:4}（对偶定理）。
\item 原貌：仅写将式 \eqref{eq:28}–\eqref{eq:30} 中每个矩阵 $\bm{X}(t_0 + \tau)$ 替换为 $\hat{\bm{X}}^{\mathrm{T}}(T - \tau)$ 即得式 \eqref{eq:35}–\eqref{eq:37}，反之亦然。
\item 订正：在完成上述逐矩阵元素替换后，补入“再对所得每一条矩阵方程两边同时取整体转置（并利用对称矩阵性质 $\hat{\bm{P}}^{\text{✫}}(t)^{\mathrm{T}} = \hat{\bm{P}}^{\text{✫}}(t)$ 与 $\hat{\bm{Q}}^{\mathrm{T}}(t) = \hat{\bm{Q}}(t)$）”。
\item 依据：以式~\eqref{eq:28} 为例，若仅执行字面替换 $\bm{X} \mapsto \hat{\bm{X}}^{\mathrm{T}}$（从而 $\bm{X}^{\mathrm{T}} \mapsto \hat{\bm{X}}$），所得等式为
\begin{equation*}
\hat{\bm{\Delta}}^{\text{✫}}(T - \tau)^{\mathrm{T}} = \hat{\bm{\Phi}}^{\mathrm{T}}(T - \tau + 1; T - \tau)\hat{\bm{P}}^{\text{✫}}(T - \tau)^{\mathrm{T}}\hat{\bm{M}}(T - \tau)\bigl[\hat{\bm{M}}^{\mathrm{T}}(T - \tau)\hat{\bm{P}}^{\text{✫}}(T - \tau)^{\mathrm{T}}\hat{\bm{M}}(T - \tau)\bigr]^{-1}
\end{equation*}
这是一个 $n \times m$ 阶矩阵方程，并非 $m \times n$ 阶的式~\eqref{eq:35}；唯有对该方程两边再同时取一次整体转置，方能逆转矩阵乘积顺序化为式~\eqref{eq:35}。同理，式~\eqref{eq:29} 与式~\eqref{eq:30} 经字面替换后亦须两边同时取整体转置，方能分别化为式~\eqref{eq:36} 中的 $-\hat{\bm{M}}\hat{\bm{\Delta}}^{\text{✫}}$ 项顺序以及式~\eqref{eq:37} 右端的 $\hat{\bm{\Phi}}^{\mathrm{T}}\hat{\bm{P}}^{\text{✫}}\hat{\bm{\Phi}}^{\text{✫}}$ 因子顺序。
\end{itemize}

\subsection{K-18：第 9 页算例 1 闭环状态转移矩阵右上元提前置零违反式 (29) 与误差方程 (23)}\label{sec:k18}

\begin{itemize}
\item 定位：第 9 页左栏算例~\ref{ex:1} 的 $\bm{\Phi}^{\text{✫}}(t + 1; t)$ 表达式及右栏讨论段。
\item 原貌：声称右上元 $(1, 2)$ 为 $0$ 的二阶矩阵 $\operatorname{diag}\bigl(\phi_{11}[1 - C(t)], 0\bigr)$ 满足差分方程~\eqref{eq:28}–\eqref{eq:30}。
\item 订正：严格满足式 \eqref{eq:28}–\eqref{eq:30} 及误差方程~\eqref{eq:23} 的矩阵订正为
\begin{equation*}
\bm{\Phi}^{\text{✫}}(t + 1; t) = \begin{bmatrix} \phi_{11}[1 - C(t)] & -\phi_{11}C(t) \\ 0 & 0 \end{bmatrix}
\end{equation*}
并在后文指明：仅在状态估计递推式~\eqref{eq:21} 中，因恒有 $x_2^{\text{✫}}(t \mid t - 1) = 0$，方可将作用在该零分量上的系数 $\phi_{12}^{\text{✫}} = -\phi_{11}C(t)$ 等效置零以实现一阶预测器。
\item 依据：算例~\ref{ex:1} 中 $\bm{\Phi}(t + 1; t) = \operatorname{diag}(\phi_{11}, 0)$，$\bm{M}(t) = \begin{bmatrix} 1 & 1 \end{bmatrix}$，$\bm{\Delta}^{\text{✫}}(t) = \begin{bmatrix} \phi_{11}C(t) & 0 \end{bmatrix}^{\mathrm{T}}$，直接代入式~\eqref{eq:29} $\bm{\Phi} - \bm{\Delta}^{\text{✫}}\bm{M}$ 立得右上元为 $-\phi_{11}C(t)$，当 $t \geq 1$（且 $a^2 > 0, \phi_{11} \neq 0$）时严格非零。更关键的是，在误差方程~\eqref{eq:23} 中，$\tilde{x}_2(t \mid t - 1) = x_2(t) = u_2(t - 1)$ 的方差为 $b^2 > 0$ 而非零；若误将 $\bm{\Phi}^{\text{✫}}$ 右上元置零并代入式~\eqref{eq:24} 的对称协方差展开 $\bm{\Phi}^{\text{✫}}\bm{P}^{\text{✫}}\bm{\Phi}^{\text{✫}}(t + 1; t)^{\mathrm{T}} + \bm{Q}$，算出的 $(1, 1)$ 元将漏掉关键噪声耦合方差项 $\phi_{11}^2 C^2(t)b^2$ 而导致计算错误。
\end{itemize}

\subsection{K-19：第 9 页算例 2 漏列初始噪声统计条件且初值预测增益矩阵第三分量算错}\label{sec:k19}

\begin{itemize}
\item 定位：第 9 页右栏，算例~\ref{ex:2} 条件 (2)–(3) 及紧随其后的 $\bm{\Delta}^{\text{✫}}(0)$ 初值式。
\item 原貌：条件中未指明初始相关噪声 $x_3(0)$ 的统计特性，并将初值一步预测增益印为 $\bm{\Delta}^{\text{✫}}(0) = \begin{bmatrix} 0 & 0 & 1 \end{bmatrix}^{\mathrm{T}}$。
\item 订正：补全初值假设 $\mathbb{E}[x_3(0)] = 0$，$\mathbb{E}[x_3^2(0)] = \sigma_3^2 > 0$ 及各初值与激励序列相互正交，并将初值预测增益订正为 $\bm{\Delta}^{\text{✫}}(0) = \begin{bmatrix} 0 & 0 & \phi_{33} \end{bmatrix}^{\mathrm{T}}$。
\item 依据：在 $t = 0$ 时 $\bm{P}^{\text{✫}}(0) = \operatorname{diag}(0, a^2, \sigma_3^2)$（须有 $\sigma_3^2 > 0$ 方能保证 $\bm{M}\bm{P}^{\text{✫}}(0)\bm{M}^{\mathrm{T}} = \sigma_3^2 > 0$ 可逆），代入式~\eqref{eq:28} 得
\begin{equation*}
\bm{\Delta}^{\text{✫}}(0) = \begin{bmatrix} 1 & 1 & 0 \\ 0 & 1 & 0 \\ 0 & 0 & \phi_{33} \end{bmatrix} \begin{bmatrix} 0 \\ 0 \\ \sigma_3^2 \end{bmatrix} \frac{1}{\sigma_3^2} = \begin{bmatrix} 0 \\ 0 \\ \phi_{33} \end{bmatrix}
\end{equation*}
物理上首拍观测为 $y_1(0) = x_1(0) + x_3(0) = x_3(0)$，故对下一拍噪声的一步预测必为 $x_3^{\text{✫}}(1 \mid 0) = \phi_{33}y_1(0)$；原刊误将当拍滤波增益 $\bm{\Phi}(0; 1)\bm{\Delta}^{\text{✫}}(0) = \begin{bmatrix} 0 & 0 & 1 \end{bmatrix}^{\mathrm{T}}$ 当成了预测增益 $\bm{\Delta}^{\text{✫}}(0)$。若将原刊错值代入式 \eqref{eq:29}–\eqref{eq:30}，所得 $\bm{P}^{\text{✫}}(1)$ 右下元将变为 $\phi_{33}(\phi_{33} - 1)\sigma_3^2 + b^2 \neq b^2$，与同式左端正确的 $\bm{P}^{\text{✫}}(1)$ 自相矛盾。
\end{itemize}

\subsection{K-20：第 10 页算例 2 当拍滤波状态递推矩阵右因子时标笔误且 (1,2) 元误印为 C\_3(t)}\label{sec:k20}

\begin{itemize}
\item 定位：第 10 页左栏上方，算例~\ref{ex:2} 的 $3 \times 3$ 当拍滤波状态转移矩阵等式。
\item 原貌：等式左端印为 $\bm{\Phi}(t; t + 1)\bm{\Phi}^{\text{✫}}(t + 1; t)\bm{\Phi}(t + 1; t)$，等式右端矩阵第一行第二列印为 $C_1(t) - tC_3(t)$（全文从未定义 $C_3$）。
\item 订正：按式~\eqref{eq:31} 取 $s = 0$ 将左端右因子时标订正为 $\bm{\Phi}(t; t - 1)$，并将右端矩阵 $(1, 2)$ 元中的 $C_3(t)$ 订正为 $C_2(t)$。
\item 依据：在式~\eqref{eq:31} 中令 $s = 0$，作用于 $\bm{x}^{\text{✫}}(t - 1 \mid t - 1)$ 的状态递推矩阵恰为 $\bm{\Phi}(t; t + 1)\bm{\Phi}^{\text{✫}}(t + 1; t)\bm{\Phi}(t; t - 1)$（当时变 $\phi_{33}(t + 1; t)$ 随时间变化时，右因子取 $\bm{\Phi}(t; t - 1)$ 才与右端第三列中的上一拍系数对应）。利用式~\eqref{eq:29} 直接展开矩阵乘积 $\bigl[\bm{I} - \bm{\Phi}(t; t + 1)\bm{\Delta}^{\text{✫}}(t)\bm{M}\bigr]\bm{\Phi}(t; t - 1)$，第一行 $\frac{1}{C_1(t)}\begin{bmatrix} C_1(t) - tC_2(t) & 0 & -tC_2(t) \end{bmatrix}$ 右乘 $\bm{\Phi}(t; t - 1)$ 的第二列 $\begin{bmatrix} 1 & 1 & 0 \end{bmatrix}^{\mathrm{T}}$，立得 $(1, 2)$ 元精确等于 $\bigl(C_1(t) - tC_2(t)\bigr) / C_1(t)$。另附备考：紧随其后的式~\eqref{eq:39} 系由连续微分近似积分所得，精确离散求和式在次主项 $O(t^2)$ 系数上相差梯形离散化余项 $\frac{1}{2}(1 - \phi_{33})^2 t^2$，但在 $t \gg 1$ 下立方主项 $\frac{1}{3}(1 - \phi_{33})^2 t^3$ 渐近精确，不影响后续全部极限结论。
\end{itemize}

\subsection{K-21：第 10 页算例 2 直接观测变量误标最优估计上标}\label{sec:k21}

\begin{itemize}
\item 定位：第 10 页左栏下半部，算例~\ref{ex:2} 渐近噪声估计式。
\item 原貌：将测量值 $y_1(t)$ 误印为带最优星号上标的估计量形式。
\item 订正：删去 $y_1(t)$ 上误加的最优上标，还原为已知观测值 $y_1(t)$，即 $x_3^{\text{✫}}(t \mid t) \approx y_1(t) - x_1^{\text{✫}}(t \mid t)$。
\item 依据：$y_1(t)$ 是直接测得的已知观测量；且因 $y_1(t) = x_1(t) + x_3(t) \in \mathcal{Y}(t)$，向 $\mathcal{Y}(t)$ 作线性正交投影实有对一切 $t \geq 0$ 严格成立的恒等式 $x_3^{\text{✫}}(t \mid t) = y_1(t) - x_1^{\text{✫}}(t \mid t)$（文中在 $t \gg 1$ 处写约等号意指位置估计 $x_1^{\text{✫}}(t \mid t)$ 已渐近收敛而几乎不再随当拍新观测波动，从而新观测的增量几乎全部归入噪声估计）。
\end{itemize}

\subsection{K-22：第 11 页附录去均值过程记号冲突及高阶混合偏导分母非标准逗号}\label{sec:k22}

\begin{itemize}
\item 定位：第 11 页左栏不相关过程中心化段落，以及右栏式~\eqref{eq:43} 下方联合概率密度定义式。
\item 原貌：左栏以右上撇号定义去均值中心化过程，与紧接式~\eqref{eq:42} 及全文用撇号表示矩阵转置的约定直接冲突；右栏高阶混合偏导数分母在各偏微分符号间印有逗号。
\item 订正：中心化过程改记为圆圈上标 $\mathring{x}(t) := x(t) - \mathbb{E}[x(t)]$；混合偏导分母去掉多余逗号写作 $\partial\xi_1 \cdots \partial\xi_n$。
\item 依据：消除同一局部上下文内单字符多义冲突，符合多元微积分标准偏导数记法。
\end{itemize}

\subsection{K-23：第 11 页附录定理 5 (A) 与 (B) 逻辑跳跃及缺失零均值与联合正态前提}\label{sec:k23}

\begin{itemize}
\item 定位：第 11 页右栏末，定理~\ref{thm:5}。
\item 原貌：(\emph{A}) 仅写“正态随机过程上的线性函数（并因此条件期望）是正态随机变量”；(\emph{B}) 写作“正交的正态随机变量相互独立”；(\emph{C}) 将一般随机过程的二阶原点混合矩 $\mathbb{E}[x(t)x(s)]$ 混称为协方差。
\item 订正：(\emph{A}) 补全“正态随机过程的任意有限维条件概率分布仍为正态分布，且条件均值是给定观测变量的仿射（零均值时为线性）函数，故条件数学期望亦为正态随机变量”；(\emph{B}) 订正为“同一正态随机过程上的不相关随机变量（或零均值正交随机变量）相互独立”；(\emph{C}) 明确区分均值函数与半正定中心协方差函数（或相容二阶混合矩）。
\item 反例与依据：针对 (\emph{A})，第 3 页推论证明引用定理~\ref{thm:5} (\emph{A}) 旨在断言“条件概率分布本身是正态分布”，而原刊定理~\ref{thm:5} (\emph{A}) 漏写了条件分布的正态性且颠倒了推导次序（须先由多元正态密度商式~\eqref{eq:44} 导出条件分布为正态且条件均值为观测的仿射组合，方能推出条件期望是仿射泛函）。针对 (\emph{B})，按同页左栏定义，“正交”指二阶原点矩 $\mathbb{E}[XY] = 0$，在非零均值下并不等价于不相关 $\operatorname{cov}(X, Y) = 0$。构造二元联合正态反例：取联合正态随机向量 $(X, Y)^{\mathrm{T}}$，令均值 $\mathbb{E}[X] = \mathbb{E}[Y] = 1$，中心协方差矩阵为正定矩阵
\begin{equation*}
\bm{C} = \begin{bmatrix} 2 & -1 \\ -1 & 2 \end{bmatrix}
\end{equation*}
此时 $\mathbb{E}[XY] = \operatorname{cov}(X, Y) + \mathbb{E}[X]\mathbb{E}[Y] = -1 + 1 = 0$，即 $X$ 与 $Y$ 严格正交且联合正态，但因 $\operatorname{cov}(X, Y) = -1 \neq 0$，二者显然不独立！此外若缺失“联合正态（属于同一正态随机过程）”前提（例如设 $X$ 服从标准正态分布，独立随机符号 $S$ 以等概率 $1/2$ 取 $\pm 1$，令 $Y = SX$，则 $X$ 与 $Y$ 各自均为零均值标准正态且 $\mathbb{E}[XY] = 0$，但 $|Y| = |X|$ 显然不独立），结论亦不成立。
\end{itemize}

\subsection{附：2002 年重排转录底本录入与光学识别错字订正汇总（L-01 至 L-31）}\label{sec:l-all}

2002 年电子重排底本系由原刊扫描图像经光学字符识别提取正文并手工重排全部公式所得，其中混入了若干非原刊固有的字形误识与公式敲击笔误，本中文版已悉数按正确文意与数学语法译出，主要条目备考如下：

\begin{itemize}
\item 正文拼写，断字与标点误植（L-01 至 L-03，L-06，L-08 至 L-12，L-15，L-17，L-20，L-21，L-23，L-24，L-26，L-29）：第 1 页文献引用\cite{r02} 右方括号误作圆括号（L-01），底注月份 \texttt{April} 末字母误识为数字 \texttt{1}（L-02），右栏末跨页句漏系动词 \texttt{is}（L-03）；第 2 页与第 9 页两处介词 \texttt{of} 误识为 \texttt{or}（L-06，L-26）；第 3 页 \texttt{Theorem 1-a} 数字 \texttt{1} 误作大写字母 \texttt{I}（L-08），跨行残留连字符 \texttt{re- quire} 与 \texttt{co-ordi- nates}（L-09）；第 4 页子空间 $\mathcal{Y}(t)$ 前后误加句点或撇号，\cite{r16} 多空格，定理~\ref{thm:2} 首句漏助动词 \texttt{be}（L-10，L-11）；第 5 页两处句末漏句号且矩阵指数与 \texttt{in} 粘连（L-12）；第 7 页重复连词 \texttt{that}（L-15），公式与 \texttt{will} 粘连（L-17），\texttt{observable} 后漏逗号（L-20）；第 8 页句中 \texttt{where} 与 \texttt{one} 首字母误大写（L-21，L-23）；第 9 页 \texttt{Results} 跨行误识为 \texttt{Resuits}（L-24）；第 11 页 \texttt{orthogonal} 前漏不定冠词 \texttt{an}（L-29）。
\item 公式符号，上下标与帽号漏误（L-04，L-05，L-07，L-13，L-14，L-16，L-18，L-19，L-21，L-22，L-24，L-25，L-29 至 L-31）：第 2 页外积定义中列向量 $\bm{x}$ 多排转置号（L-04），紧接外积期望 $\mathbb{E}[\bm{x}\bm{y}^{\mathrm{T}}]$ 将粗体 $\bm{y}$ 误识为粗体 $\bm{v}$（L-05）；第 3 页式~\eqref{eq:3} 右端 $x_1(t_1)$ 漏下标 \texttt{1}（L-07）；第 6 页首行稳态级数左端向量 $\bm{x}(t)$ 漏加粗排成正体（L-13）；第 7 页单步余投影左端多排右圆括号且花体 $\mathcal{Z}(t)$ 误排斜体（L-14），协方差引用 $\bm{P}^{\text{✫}}(t)$ 漏最优上标（L-16），多步展开求和下限 \texttt{r = 1} 漏等号排成 \texttt{r1}（L-18），式~\eqref{eq:31} 首行 $\bm{x}^{\text{✫}}(t + 1 \mid t)$ 条件竖线前多排右圆括号（L-19）；第 8 页问题~\ref{prob:II} 性能指标中 $\hat{\bm{Q}}(\tau)$ 与最优指标中 $\hat{\bm{P}}^{\text{✫}}(t - 1)$ 漏敲对偶帽号（L-21，L-22）；第 9 页评注~\ref{rem:o} 问题~\ref{prob:I} 输出矩阵 $\bm{M}(t)$ 误加帽号（L-24），算例~\ref{ex:1} 条件 (3) 中 $\mathbb{E}[u_2^2(t)] = b^2$ 漏下标 \texttt{2}（L-25）；第 11 页式~\eqref{eq:42} 下方向量 $\bm{x}$ 与多元正态密度下标随机向量 $\bm{x}$ 漏加粗（L-29，L-31），联合密度分子偏导下标首项 $x(t_1)$ 误敲为 $n(t_1)$（L-30）。
\item 参考文献书目光学识别错字（L-27，L-28）：文献\cite{r03} 刊名 \texttt{Proceedings} 误作 \texttt{Proceedinqs}；文献\cite{r09} 刊名 \texttt{Group} 误加句点作 \texttt{Grou.p}；文献\cite{r11} 题名 \texttt{Time-Varying} 误作 \texttt{Tirne-Varying}；文献\cite{r17} 作者 \texttt{O. A. Gross} 与文献\cite{r28} 作者 \texttt{O. J. M. Smith} 首字母 \texttt{O.} 误识为数字零 \texttt{0.}；文献\cite{r19} 与\cite{r20} 刊名缩写 \texttt{AIEE} 误将大写 \texttt{I} 排作小写 \texttt{l}；文献\cite{r07}，\cite{r08}，\cite{r10}，\cite{r11}，\cite{r12}，\cite{r14}，\cite{r20} 中的书目分隔标点与空格讹误亦已一并规范订正。
\end{itemize}
